\textbf{Proposición.} Sea $X$ un espacio topológico conexo y $Y$ un conjunto cualquiera. Si $f: X \\to Y$ es una función localmente constante, entonces $f$ es constante en todo $X$. \\begin{proof} Fijemos un punto $x\_0 \\in X$ y sea $c = f(x\_0) \\in Y$. Definimos los subconjuntos \\\\[ A = \\{x \\in X \\mid f(x) = c\\} \\quad \\text{y} \\quad B = \\{x \\in X \\mid f(x) \\neq c\\}. \\\\] Veamos que $A$ y $B$ son subconjuntos abiertos en $X$: \\begin{enumerate} \\item \\textbf{$A$ es abierto en $X$:} Sea $x \\in A$. Como $f$ es localmente constante, existe una vecindad abierta $U\_x$ de $x$ tal que $f$ toma un único valor en $U\_x$. Dado que $x \\in A$, se tiene $f(x) = c$, de modo que $f(y) = c$ para todo $y \\in U\_x$. Por tanto $U\_x \\subseteq A$, lo que demuestra que $A$ es abierto. \\item \\textbf{$B$ es abierto en $X$:} Sea $x \\in B$. Entonces $f(x) = d$ con $d \\neq c$. Por ser $f$ localmente constante, existe una vecindad abierta $V\_x$ de $x$ tal que $f(y) = d \\neq c$ para todo $y \\in V\_x$. Por tanto $V\_x \\subseteq B$, lo que demuestra que $B$ es abierto. \\end{enumerate} Por definición de $A$ y $B$, es claro que $A \\cap B = \\emptyset$ y $A \\cup B = X$. Además, como $x\_0 \\in A$, se tiene $A \\neq \\emptyset$. Si supusiéramos que $B \\neq \\emptyset$, el par $(A, B)$ constituiría una separación de $X$ por abiertos disjuntos no vacíos, lo cual contradice la hipótesis de que $X$ es conexo. Por lo tanto, $B = \\emptyset$, de donde $A = X$. Concluimos que $f(x) = c$ para todo $x \\in X$, es decir, $f$ es globalmente constante. \\end{proof}
